\section{Hardware Setup}
\begin{itemize}
	\item \textbf{CPU:} 12th Gen Intel(R) Core(TM) i7-12650H
	\item \textbf{RAM:} 16 GB
	\item \textbf{GPU:} NVIDIA GeForce RTX 3050 Ti Laptop GPU
\end{itemize}

\section{Software Frameworks}
To ensure the full reproducibility of the experiments, the software environment was managed using virtual environments. The core frameworks and primary dependencies utilized in this research are detailed below:

\begin{itemize}
    \item \textbf{Federated Learning Framework:} Flower (\texttt{flwr}) v1.7.0.
    \item \textbf{Deep Learning Framework:} PyTorch v2.5.1 (with CUDA 12.1 support).
    \item \textbf{Experiment Configuration:} Hydra-Core v1.3.2.
\end{itemize}

The exhaustive list of all dependencies, environment configuration files, installation instructions, and execution scripts is publicly available in the official project repository:
\begin{center}
    \url{https://github.com/IkerMardure/tfg-gossip-learning-backdoor}
\end{center}
In this repository, the \texttt{README.md} file provides the exact commands required to replicate the development environment and run the simulations presented in this thesis.

\section{Simulation Runtime}
The experimentation involves computationally intensive simulations of 240 communication
rounds for both MNIST and NEU-DET scenarios. Runtime is dominated by
repeated local training, peer-to-peer aggregation, and metric logging at each round, with
additional overhead when adversarial evaluation is enabled. On the hardware described
above, the full 240-round in the fully connected runs required approximately 24000 seconds
for MNIST and 4000 seconds for NEU-DET. While the NEU-DET scenario utilized
a significantly heavier architecture (ResNet-18), its total runtime was substantially lower
due to the massive difference in dataset volume (1,800 total images compared to the 60,000
images processed per epoch in the MNIST baseline).